% ============================================================
%  Tarea-Examen 3. Complejidad Computacional.
%  Tema: SAT y 3SAT son NP-completos (Teorema de Cook-Levin)
% ============================================================
\documentclass[12pt, letterpaper]{article}

\usepackage[utf8]{inputenc}
\usepackage[T1]{fontenc}
\usepackage[spanish, mexico]{babel}
\usepackage{amsmath, amssymb, amsthm}
\usepackage{geometry}
\usepackage{enumitem}
\usepackage{mathtools}
\usepackage{xcolor}
\usepackage{mdframed}

\geometry{margin=2.5cm}

% ── Macros ───────────────────────────────────────────────────────────────────
\newcommand{\SAT}{\textsc{Sat}}
\newcommand{\tSAT}{\textsc{3Sat}}
\newcommand{\NP}{\mathbf{NP}}
\newcommand{\PP}{\mathbf{P}}
\newcommand{\coNP}{\mathbf{coNP}}
\newcommand{\ntime}{\textsc{NTime}}
\newcommand{\leqp}{\leq_p}
\newcommand{\abs}[1]{\left|#1\right|}

% Encabezado de problema con puntaje
\newcommand{\problema}[2]{%
  \bigskip
  \noindent\textbf{Problema #1} \textnormal{(#2 puntos)}\par\smallskip
}

% ── Documento ────────────────────────────────────────────────────────────────
\begin{document}

\begin{center}
  {\large \textbf{Tarea-examen 3. Complejidad Computacional.}}\\[4pt]
  \textit{El Teorema de Cook-Levin: \SAT{} y \tSAT{} son NP-completos}
\end{center}

\smallskip
\begin{mdframed}
\textbf{Indicaciones.} Justifica todas las respuestas. Cuando demuestres una
implicación, indica explícitamente la dirección que estás demostrando.

\textbf{Los Problemas 3, 4 y 5 suman 10 puntos y son suficientes para obtener
10.}  Los Problemas 1, 2 y 6 son opcionales y sirven como puntos de
recuperación: si alguno de los problemas principales no te sale completo,
puedes compensar resolviendo alguno de estos.  El total disponible es de
\textbf{17 puntos}.

\medskip
\textbf{Entrega.}  El documento se entrega de forma \textbf{digital} a más
tardar el \textbf{jueves 11 de junio}.  Puedes escribir tus respuestas
\textbf{a mano} (foto o escaneo legible) \textbf{o en \LaTeX{}}.  En cualquier
caso, sube tu archivo a \textbf{Google Drive} o \textbf{GitHub} y comparte el
enlace (con permisos de lectura) en la asignación de Classroom.
\end{mdframed}

% ====================================================================
\problema{1}{2} \textbf{Consecuencia directa: $\PP = \NP$ si y solo si
$\SAT \in \PP$.}

En este problema formalizamos la razón por la que $\SAT$ es el problema
``más difícil'' de $\NP$.

\begin{enumerate}[label=\alph*)]

\item (0.5 pts) ($\Rightarrow$) Supón que $\PP = \NP$.  Concluye que
  $\SAT \in \PP$ y justifica brevemente por qué.

\item (1 pt) ($\Leftarrow$) Supón ahora que $\SAT \in \PP$.  Sea
  $L \in \NP$ un lenguaje arbitrario.
  \begin{enumerate}[label=\roman*.]
    \item Explica por qué existe una reducción de Karp en tiempo polinomial
      $f$ tal que $x \in L \iff f(x) \in \SAT$.  \textit{(Basta citar el
      resultado correspondiente visto en clase.)}
    \item Dado un algoritmo polinomial $D$ para $\SAT$, describe un
      algoritmo polinomial para $L$ usando $f$ y $D$.
    \item Concluye que $L \in \PP$.
  \end{enumerate}

\item (0.5 pts) Combina los incisos anteriores para enunciar la equivalencia
  $\PP = \NP \iff \SAT \in \PP$.  ¿Por qué este resultado convierte a $\SAT$
  en un problema ``representante'' de toda la clase $\NP$?

\end{enumerate}

% ====================================================================
\problema{2}{2} \textbf{Expresividad de fórmulas CNF.}

Recordamos las siguientes definiciones.  Un \emph{literal} es una variable
booleana $x_i$ o su negación $\bar{x}_i$.  Una \emph{cláusula} es una
disyunción de literales, es decir, una expresión de la forma
$(\ell_1 \vee \ell_2 \vee \cdots \vee \ell_k)$ donde cada $\ell_j$ es un
literal.  Una fórmula en \emph{Forma Normal Conjuntiva} (CNF) es una
conjunción de cláusulas:
\[
  \varphi = \bigwedge_{i=1}^{m} \left(\bigvee_{j} \ell_{ij}\right).
\]
Antes de probar que $\SAT$ es NP-completo necesitamos saber ``hablar en CNF''.

\begin{enumerate}[label=\alph*)]

\item (0.5 pts) La fórmula $(x_1 \vee \bar{y}_1) \wedge (\bar{x}_1 \vee y_1)$
  está en forma CNF.
  \begin{enumerate}[label=\roman*.]
    \item Muestra que es satisfecha \emph{si y solo si} $x_1 = y_1$.
    \item Usando este hecho, escribe una fórmula CNF $\varphi_{=}$ sobre las
      variables $x_1,\ldots,x_n,y_1,\ldots,y_n$ que es satisfecha si y solo
      si $x_i = y_i$ para todo $i \in \{1,\ldots,n\}$.  ¿Cuántas cláusulas
      tiene $\varphi_{=}$?
  \end{enumerate}

\item (1 pt) \textbf{Universalidad de AND/OR/NOT.}
  Sea $f : \{0,1\}^\ell \to \{0,1\}$ una función booleana cualquiera, con
  $\ell \in \mathbb{N}$ fijo.
  \begin{enumerate}[label=\roman*.]
    \item Para cada $v \in \{0,1\}^\ell$ tal que $f(v) = 0$, construye una
      cláusula $C_v$ en las $\ell$ variables tal que $C_v(v) = 0$ y
      $C_v(u) = 1$ para todo $u \neq v$.

      \textit{Sugerencia:} Si $v = (v_1, \ldots, v_\ell)$, considera el
      literal $z_j = x_j$ si $v_j = 0$ y $z_j = \bar{x}_j$ si $v_j = 1$.

    \item Define $\varphi_f = \bigwedge_{v:\, f(v)=0} C_v$.  Demuestra que
      $\varphi_f(u) = f(u)$ para todo $u \in \{0,1\}^\ell$, y acota el
      tamaño de $\varphi_f$.
  \end{enumerate}

\item (0.5 pts) Sea $F : \{0,1\}^k \to \{0,1\}^c$ una función con $k$ y $c$
  constantes (que dependen de la máquina, pero no del tamaño de la entrada).
  \begin{enumerate}[label=\roman*.]
    \item Explica por qué cada bit de salida de $F$ puede expresarse como una
      fórmula CNF de tamaño \emph{constante} (usando el inciso b).
    \item ¿Por qué este hecho es crucial en la prueba del Teorema de
      Cook-Levin?
  \end{enumerate}

\end{enumerate}

% ====================================================================
\problema{3}{5} \textbf{$\SAT \leq_p \tSAT$: la reducción cláusula por
cláusula.}

Recordamos que $k$-\SAT{} es el lenguaje de las fórmulas CNF satisfacibles en
las que cada cláusula tiene \emph{a lo más} $k$ literales.  El objetivo de este
problema es mostrar que toda fórmula CNF puede convertirse en una fórmula 3CNF
equivalente (en cuanto a satisfacibilidad) en tiempo polinomial.

\begin{enumerate}[label=\alph*)]

\item (1 pt) \textbf{Cláusulas cortas.}
  \begin{enumerate}[label=\roman*.]
    \item Sea $C = (\ell_1)$ una cláusula de un solo literal.  Introduce dos
      variables auxiliares $z_1, z_2$ y reemplaza $C$ por una conjunción de
      cláusulas de exactamente 3 literales, preservando la satisfacibilidad.
    \item Sea $C = (\ell_1 \vee \ell_2)$ una cláusula de dos literales.
      Introduce una variable auxiliar $z$ y reemplaza $C$ por una cláusula de
      exactamente 3 literales, preservando la satisfacibilidad.
  \end{enumerate}

\item (1 pt) \textbf{Cláusula de 4 literales.}
  Sea $C = (\ell_1 \vee \ell_2 \vee \ell_3 \vee \ell_4)$.  Introduce una
  variable auxiliar $z$ y propón dos cláusulas $C_1$ y $C_2$ de exactamente 3
  literales que reemplacen a $C$.

\item (1.5 pts) \textbf{Correctud del reemplazo.}
  Demuestra que $C_1 \wedge C_2$ (del inciso b) es satisfacible si y solo si
  $C$ lo es.
  \begin{enumerate}[label=\roman*.]
    \item ($\Rightarrow$) Si existe una asignación que satisface $C$, muestra
      que existe un valor para $z$ que satisface $C_1 \wedge C_2$.
    \item ($\Leftarrow$) Si $C_1 \wedge C_2$ es satisfacible, muestra que
      la misma asignación (ignorando $z$) satisface $C$.
  \end{enumerate}

\item (1 pt) \textbf{Generalización: cláusula de tamaño $k > 3$.}
  Sea $C = (\ell_1 \vee \ell_2 \vee \cdots \vee \ell_k)$ con $k > 3$.
  \begin{enumerate}[label=\roman*.]
    \item Describe una construcción recursiva que reemplaza $C$ por una
      conjunción de cláusulas de tamaño $\leq 3$, usando variables auxiliares
      $z_1, \ldots, z_{k-3}$.

      \textit{Sugerencia:} En el paso base, aplica el inciso (b) con
      $\ell_1, \ell_2$ y $z_1$ para obtener una cláusula $(\ell_1 \vee \ell_2
      \vee z_1)$ y luego recurre sobre $(\bar{z}_1 \vee \ell_3 \vee \cdots
      \vee \ell_k)$.

    \item ¿Cuántas variables auxiliares se introducen?
    \item ¿Cuántas cláusulas de tamaño 3 produce la construcción?
    \item Argumenta (por inducción o por el mismo razonamiento del inciso c)
      que la conjunción resultante es satisfacible si y solo si $C$ lo es.
  \end{enumerate}

\item (0.5 pts) \textbf{La reducción $f$.}
  Sea $\varphi$ una fórmula CNF con $m$ cláusulas y longitud total
  $s = |\varphi|$ (número de símbolos).  Aplica la construcción anterior
  cláusula por cláusula para obtener una fórmula 3CNF $\psi = f(\varphi)$.
  \begin{enumerate}[label=\roman*.]
    \item Acota el tamaño de $\psi$ en función de $s$.
    \item Muestra que $f$ es computable en tiempo polinomial.
    \item Concluye que $\SAT \leq_p \tSAT$.
  \end{enumerate}

\end{enumerate}

% ====================================================================
\problema{4}{2} \textbf{3SAT es NP-completo.}

Usa los resultados de los problemas anteriores para concluir el Teorema de
Cook-Levin.

\begin{enumerate}[label=\alph*)]

\item (0.5 pts) Muestra que $\tSAT \in \NP$ describiendo un testigo y un
  verificador polinomial.

\item (1 pt) Usa que $\SAT$ es NP-duro (visto en clase) y la reducción
  $\SAT \leq_p \tSAT$ (Problema 3e) para mostrar que $\tSAT$ es NP-duro.
  \textit{Sugerencia:} usa la transitividad de las reducciones en tiempo
  polinomial.

\item (0.5 pts) Concluye que $\tSAT$ es NP-completo.

\end{enumerate}

% ====================================================================
\problema{5}{3} \textbf{Autoreducibilidad de \SAT: decidir implica
encontrar.}

En este problema mostramos que si $\PP = \NP$ (y por lo tanto $\SAT \in \PP$),
entonces no solo podemos \emph{decidir} si una fórmula es satisfacible, sino
también \emph{encontrar} una asignación satisfactoria en tiempo polinomial.

Sea $D$ un algoritmo en tiempo polinomial para $\SAT$ y sea $\varphi$ una
fórmula CNF satisfacible sobre las variables $x_1, \ldots, x_n$.  Para
$b \in \{0,1\}$, denotamos por $\varphi[x_i \leftarrow b]$ la fórmula
obtenida al sustituir $x_i$ por $b$ y simplificar.

\begin{enumerate}[label=\alph*)]

\item (0.5 pts) Muestra que si $\varphi \in \SAT$, entonces
  $\varphi[x_1 \leftarrow 0] \in \SAT$ o
  $\varphi[x_1 \leftarrow 1] \in \SAT$ (o ambas).

\item (1 pt) Describe el siguiente algoritmo en pseudocódigo y argumenta
  que es correcto:
  \begin{enumerate}[label=\roman*.]
    \item Para $i = 1, \ldots, n$: si $D(\varphi[x_i \leftarrow 0]) = 1$,
      fija $x_i = 0$; de lo contrario, fija $x_i = 1$.  Actualiza $\varphi$
      sustituyendo el valor asignado a $x_i$.
    \item Al terminar el ciclo, devuelve la asignación $(x_1, \ldots, x_n)$.
  \end{enumerate}
  Muestra que, si $\varphi \in \SAT$, el algoritmo devuelve una asignación
  que satisface $\varphi$.

\item (0.5 pts) ¿Cuántas llamadas a $D$ realiza el algoritmo?
  Concluye que el algoritmo corre en tiempo polinomial (suponiendo que $D$
  es polinomial).

\item (1 pt) \textbf{Extensión a cualquier $L \in \NP$.}
  Sea $L \in \NP$ con reducción de Cook-Levin $x \mapsto \varphi_x$
  (vista en clase), donde la asignación satisfactoria de $\varphi_x$ incluye,
  entre sus variables, al testigo $u$ de $x \in L$.
  Suponiendo $\PP = \NP$, describe cómo encontrar un testigo para $x \in L$
  en tiempo polinomial.

\end{enumerate}

% ====================================================================
\problema{6}{3} \textbf{$\tSAT \leq_p \textsc{IndSet}$: las reducciones
se propagan.}

Definimos
$\textsc{IndSet} = \{\langle G, k\rangle :\text{$G$ tiene un conjunto
independiente de tamaño $k$}\}$.
En este problema usamos que $\tSAT$ es NP-completo (Problema 4) para mostrar
que $\textsc{IndSet}$ también lo es.

Sea $\varphi = C_1 \wedge \cdots \wedge C_m$ una fórmula 3CNF con $m$
cláusulas.  Para cada cláusula $C_i = (\ell_{i,1} \vee \ell_{i,2} \vee
\ell_{i,3})$ asociamos un \emph{grupo} de tres vértices
$\{v_{i,1}, v_{i,2}, v_{i,3}\}$ en una gráfica $G(\varphi)$.  Las aristas
de $G(\varphi)$ se agregan según las siguientes dos reglas:
\begin{enumerate}[label=(R\arabic*)]
  \item \textbf{Aristas internas.}  Para cada grupo $i$ y para todo par
    $j \neq j'$, agrega la arista $\{v_{i,j}, v_{i,j'}\}$.
  \item \textbf{Aristas de conflicto.}  Para grupos distintos $i \neq i'$
    y cualesquiera $j, j'$, agrega la arista $\{v_{i,j}, v_{i',j'}\}$ si
    $\ell_{i,j}$ y $\ell_{i',j'}$ son literales contradictorios (i.e.
    $\ell_{i,j} = x_k$ y $\ell_{i',j'} = \bar{x}_k$ para alguna variable
    $x_k$, o viceversa).
\end{enumerate}

\begin{enumerate}[label=\alph*)]

\item (0.5 pts) Da un ejemplo pequeño: toma $\varphi = (x_1 \vee x_2 \vee
  \bar{x}_3) \wedge (\bar{x}_1 \vee x_2 \vee x_3)$, dibuja $G(\varphi)$
  indicando vértices, aristas internas y aristas de conflicto, y señala un
  conjunto independiente de tamaño $m = 2$.

\item (1 pt) ($\varphi$ satisfacible $\Rightarrow$ $G(\varphi)$ tiene
  un conjunto independiente de tamaño $m$.)  Sea $\sigma$ una asignación
  que satisface $\varphi$.
  \begin{enumerate}[label=\roman*.]
    \item Para cada cláusula $C_i$, elige un literal $\ell_{i,j_i}$ que
      sea verdadero bajo $\sigma$ y sea el vértice representante de $C_i$.
      Define $S = \{v_{i,j_i} : 1 \leq i \leq m\}$.
    \item Muestra que $|S| = m$.
    \item Muestra que $S$ es un conjunto independiente (usa las reglas R1
      y R2 para descartar que haya aristas dentro de $S$).
  \end{enumerate}

\item (1 pt) ($G(\varphi)$ tiene un conjunto independiente de tamaño $m$
  $\Rightarrow$ $\varphi$ es satisfacible.)  Sea $S$ un conjunto independiente
  de tamaño $m$ en $G(\varphi)$.
  \begin{enumerate}[label=\roman*.]
    \item Usa la regla R1 para mostrar que $S$ contiene exactamente un
      vértice de cada grupo.
    \item Usa la regla R2 para mostrar que los literales correspondientes
      a los vértices de $S$ son pairwise consistentes (nunca aparecen $x_k$
      y $\bar{x}_k$ simultáneamente).
    \item Construye una asignación $\sigma$ que satisfaga $\varphi$.
  \end{enumerate}

\item (0.5 pts) Muestra que la construcción $\varphi \mapsto
  \langle G(\varphi), m \rangle$ es computable en tiempo polinomial en
  $|\varphi|$.  Concluye que $\tSAT \leq_p \textsc{IndSet}$ y, usando que
  $\tSAT$ es NP-duro y que $\textsc{IndSet} \in \NP$, que
  $\textsc{IndSet}$ es NP-completo.

\end{enumerate}

\end{document}
